% ------------------------------------------------------------------------
%
\begin{frame}{Boolean expressions}

  Let~\(\wedge\) represent the conjunction (`and'), \(\vee\) the
  disjunction (`or'), \(\neg\) the negation (`not'), and let the
  two values \true and \false.

  \bigskip

  The definition of the conjunction of values is very simple:
  \begin{mathpar}
    \inferrule
       {}
       {\true \wedge \true \rightarrow \true}
    \and
    \inferrule
       {}
       {\true \wedge \false \rightarrow \false}\\
    \inferrule
       {}
       {\false \wedge \true \rightarrow \false}
    \and
    \inferrule
       {}
       {\false \wedge \false \rightarrow \false}
\end{mathpar}

\end{frame}

% ------------------------------------------------------------------------
%
\begin{frame}{Boolean expressions made simpler}

  This system can even be made simpler, by means of meta-variables:
  \begin{mathpar}
    \inferrule*[right=\(\quad\TirName{True}\)]
       {}
       {\true \wedge \true \rightarrow \true}\\
    \inferrule*[right=\(\quad\TirName{False}_1\)]
       {}
       {\false \wedge v \rightarrow \false}\\
    \inferrule*[right=\(\quad\TirName{False}_2\)]
       {}
       {v \wedge \false \rightarrow \false}
  \end{mathpar}
  where \(v\) denotes any value, that is, either \true or \false.

\end{frame}

% ------------------------------------------------------------------------
%
\begin{frame}{Boolean expressions}

  So, if we define the rewrite of boolean expressions (not just the
  values \true and \false) as
  \begin{mathpar}
    \inferrule*[right=\(\TirName{\(\wedge_1\)}\)]
       {e_1 \rightarrow e'_1}
       {e_1 \wedge\, e_2 \rightarrow e'_1 \wedge\, e_2}
    \and
    \inferrule*[right=\(\TirName{\(\wedge_2\)}\)]
       {e_2 \rightarrow e'_2}
       {e_1 \wedge\, e_2 \rightarrow e_1 \wedge\, e'_2}
  \end{mathpar}
  then, for example,
  \begin{mathpar}
    \inferrule*[right=\(\sigma_1(\wedge_1)\)]
       {\false \wedge\; \true \rightarrow \false
        \quad \TirName{False}_1}
       {(\false \wedge\; \true) \wedge\, (\true \wedge\, \true)
         \rightarrow \false \wedge\, (\true \wedge\, \true)}
  \end{mathpar}
  where \(\sigma_1=\{e_1 \mapsto \false \wedge \true; e_2 \mapsto
  \true \wedge \true; e'_1 \mapsto \false\}\).

\end{frame}

% ------------------------------------------------------------------------
%
\begin{frame}{Boolean expressions}

  Then
  \begin{mathpar}
    \inferrule*[right=\(\sigma_2(\wedge_2)\)]
       {\true \wedge\, \true \rightarrow \true \quad \TirName{True}}
       {\false \wedge\, (\true \wedge\, \true) \rightarrow \false \wedge\,
  \true}
  \end{mathpar}
  where \(\sigma_2 = \{e_1 \mapsto \false; e_2 \mapsto \true \wedge
  \true; e'_1 \mapsto \true\}\).

  \bigskip

  and, finally, we can instantiate \(\TirName{False}_1\) with
  \(\sigma_3 = \{v \mapsto \true\}\):
  \begin{mathpar}
    \inferrule*[right=\(\quad\sigma_3(\TirName{False}_1)\)]
       {}
       {\false \wedge\, \true \rightarrow \false}
  \end{mathpar}
  In this example, we computed both arguments of \(\wedge\) to reach
  the value.

\end{frame}

% ------------------------------------------------------------------------
%
\begin{frame}{Boolean expressions}

  But the inspection of the basic rules shows that if one of the
  arguments of \(\wedge\) reduces to \false, then there is no need to
  reduce the other argument: the rewrite will always be
  \false. Therefore, it is more efficient to extend the rules on
  values to the expressions:
  \begin{mathpar}
    \inferrule*[right=\(\quad\TirName{False}_1\)]
       {}
       {\false \wedge e \rightarrow \false}
    \and
    \inferrule*[right=\(\quad\TirName{False}_2\)]
       {}
       {e \wedge \false \rightarrow \false}\\
    \inferrule*[right=\(\quad\TirName{True}_1\)]
       {}
       {\true \wedge e \rightarrow e}
    \and
    \inferrule*[right=\(\quad\TirName{True}_2\)]
       {}
       {e \wedge \true \rightarrow e}
  \end{mathpar}
  where \(e\) represents any boolean expression. Then, now
  \begin{mathpar}
    \inferrule*[right=\(\quad\sigma_3(\TirName{False}_1)\)]
       {}
       {\false \wedge\, (\true \wedge\, \true) \rightarrow \false}
  \end{mathpar}
  where \(\sigma_3 = \{e \mapsto \true \wedge \true\}\).

\end{frame}

% ------------------------------------------------------------------------
%
\begin{frame}[fragile]{Boolean expressions}

  Try to reduce
  \begin{equation*}
    (\true \wedge (\false \wedge \true)) \wedge (\true \wedge (\true
    \wedge \true))
  \end{equation*}
  Note that this rewrite system is not deterministic. If we enforce
  that the left argument of \(\wedge\) is completely reduced first,
  then we can check if it is \true or \false: in the latter case, we
  do not need to reduce the right argument.

  \begin{mathpar}
    \inferrule*[right=\(\quad\wedge_{\TirName{True}}\)]
       {}
       {\true \wedge e \rightarrow e}
    \and
    \inferrule*[right=\(\quad\wedge_{\TirName{False}}\)]
       {}
       {\false \wedge e \rightarrow \false}\\
    \inferrule*[right=\(\wedge\)]
       {e_1 \rightarrow e'_1}
       {e_1 \wedge\, e_2 \rightarrow e'_1 \wedge\, e_2}
  \end{mathpar}
  Many programming languages, as C, use this strategy. It allows the
  programmers to write handy code like
{\footnotesize
\begin{verbatim}
while (index < max && tab[index] > 0) ...
\end{verbatim}
}

\end{frame}

% ------------------------------------------------------------------------
%
\begin{frame}{Boolean expressions}

  Note that this strategy (left argument before the right one) is not
  optimal in general: the first argument may reduce to
  \true. \textbf{Finding an optimal strategy is not always possible.}
  For example, we could add another rule that shortens sometimes the
  reductions:
  \begin{mathpar}
    \inferrule*[right=\(\wedge_{\TirName{Refl}}\)]
       {}
       {e \wedge e \rightarrow e}
  \end{mathpar}
  This rule explicitly specifies the \textbf{reflexivity} of the
  conjunction.

  But the system becomes non-deterministic because there is now an
  overlap between rule \(\wedge_{\TirName{Refl}}\) and both
  \(\wedge_{\proc{True}}\) and \(\wedge_{\proc{False}}\) (consider for
  example \(\true \wedge \false\)).

\end{frame}

% ------------------------------------------------------------------------
%
\begin{frame}{Determinism versus non-determinism}

  Let us summarise the concepts of determinism and
  non\hyp{}determinism.
  \begin{itemize}

    \item Determinism means that, for any expression, there is only
      one rule to rewrite it. This implies that, every time a program
      is run with the same input, the output will be the same. This is
      a useful property, called soundness, in particular for
      debugging.

    \item Non-determinism means that, for any expression, there may be
      several rules to rewrite it. Moreover,
      \begin{itemize}

         \item if the rewrite system is sound, the output will always
           be the same for any run on the same input, even if the
           trace may be different each time, which makes the debugging
           more difficult, but enable parallel reductions.

         \item the rewrite system may be unsound.

      \end{itemize}
  \end{itemize}

\end{frame}

% ------------------------------------------------------------------------
%
\begin{frame}[fragile]{Determinism versus non-determinism}

  Some programming languages have features that may be both
  non\hyp{}deterministic and unsound.

  \bigskip

  Perhaps the most infamous case is in C, where variables do not need
  to be initialised at their declaration:
{\footnotesize
\begin{verbatim}
int a;
\end{verbatim}
}
  If the programmer forgets to give a value to the variable, any access
  to it is unspecified.

  \bigskip

  Practically, this means that whatever is located in the memory at
  the location of the variable will be used at each read access. Of
  course, this value may change from one run of the program to
  another, even on the same input.

\end{frame}

% ------------------------------------------------------------------------
%
\begin{frame}{Commutativity and its consequences}

  We may try another rewrite system for \(\wedge\):
  \begin{mathpar}
    \inferrule*[right=\(\quad\TirName{True}\)]
       {\true \wedge e \rightarrow e}
       {}
    \and
    \inferrule*[right=\(\quad\TirName{False}\)]
       {}
       {\false \wedge e \rightarrow \false}\\
    \inferrule*[right=\(\wedge\)]
       {e_1 \rightarrow e'_1}
       {e_1 \wedge e_2 \rightarrow e'_1 \wedge e_2}
    \and
    \inferrule*[right=\(\TirName{Comm}\)]
       {}
       {e_1 \wedge e_2 \rightarrow e_2 \wedge e_1}
  \end{mathpar}
  While this system is in theory fine, there are two practical
  problems:
  \begin{enumerate}

    \item it is \textbf{not deterministic} due to \(\TirName{Comm}\)
      that overlaps with all the other rules;

    \item it may \textbf{not terminate} due to rule
      \(\TirName{Comm}\):
      \begin{align*}
        \true \wedge \false
        & \;\xrightarrow{\TirName{Comm}}\;
        \false \wedge \true \;\xrightarrow{\TirName{Comm}}\; \true
        \wedge \false\\
        & \;\xrightarrow{\TirName{Comm}}\; \false
        \wedge \true \rightarrow \cdots
      \end{align*}

  \end{enumerate}

\end{frame}

% ------------------------------------------------------------------------
%
\begin{frame}[fragile]{Non-termination}

  Now some reductions may not be finite, which models
  \textbf{non\hyp{}termination}.

  \bigskip

  This particular kind of non\hyp{}termination, due to commutativity,
  is the model in programming languages of so\hyp{}called
  `\textbf{infinite loops}' that do not allocate memory. Equivalently,
  it corresponds to a \emph{recursive call} like
{\footnotesize
\begin{verbatim}
bool and(bool e1, bool e2) { ... return and(e2, e1); ... }
\end{verbatim}
}
  Anyway, in the presence of infinite reductions, the soundness
  property must be weakened as follows:

  \bigskip

  \begin{quote}
    \textbf{Weak soundness.} \emph{All finite reductions of an
    expression end with the same value, if any.}
  \end{quote}

\end{frame}

% ------------------------------------------------------------------------
%
\begin{frame}{Non-termination}

  Let us define the negation:
  \begin{mathpar}
    \inferrule
       {}
       {\neg \true \rightarrow \false}
    \and
    \inferrule
       {}
       {\neg \false \rightarrow \true}
  \end{mathpar}
  It is also well known that the negation satisfies \(e = \neg\neg
  e\). So we could imagine another rewrite rule
  \begin{equation*}
    \neg\neg e \rightarrow e
  \end{equation*}
  This is fine, this rule simplifies double negations. But what if we
  had the converse rule?
  \begin{equation*}
    e \rightarrow \neg\neg e \qquad \text{where} \; e \; \text{is not a
  value.}
  \end{equation*}
  This allows infinite reductions: \(e \rightarrow \neg\neg e
  \rightarrow e \rightarrow \neg\neg e \rightarrow \dots\)

% excluded middle: e or ~e
% non-contradiction: ~(e and ~e)
% de Morgan applied to non-contradiction: ~(e and ~e) = ~e or ~~e

\end{frame}

% ------------------------------------------------------------------------
%
\begin{frame}{Summary}

  There are three kinds of reductions:
  \begin{enumerate}

    \item Finite and ending in a value
  \[e_1 \rightarrow e_2 \rightarrow \dots \rightarrow v\]

    \item Finite but ending in a stuck expression
  \[e_1 \rightarrow e_2 \rightarrow \dots \rightarrow e_n
    \nrightarrow\]

    \item Infinite
  \[e_1 \rightarrow e_2 \rightarrow \cdots\]

  \end{enumerate}

\end{frame}

% ------------------------------------------------------------------------
%
\begin{frame}{Conditional expressions}

  At this point, boolean expressions can be
  \begin{itemize}

    \item \Xtrue (a value)

    \item \Xfalse (a value)

    \item \(e_1 \wedge e_2\)

    \item \(e_1 \vee e_2\)

    \item \(\neg e_1\)

  \end{itemize}
  where \(e_1\) and \(e_2\) are boolean expressions. Let us extend our
  boolean expressions with a \textbf{conditional expression}:
  \begin{itemize}

     \item \Xif \(e_1\) \Xthen \(e_2\) \Xelse \(e_3\)

  \end{itemize}

\end{frame}

% ------------------------------------------------------------------------
%
\begin{frame}{Conditional expressions}

  It is easy to extend the rewrite system to cope with conditionals:
  \begin{mathpar}
    \inferrule*[right=\(\quad\TirName{If-True}\)]
       {}
       {\Xif \; \Xtrue \; \Xthen \; e_2 \; \Xelse \; e_3 \rightarrow e_2}
    \and
    \inferrule*[right=\(\quad\TirName{If-False}\)]
       {}
       {\Xif \; \Xfalse \; \Xthen \; e_2 \; \Xelse \; e_3 \rightarrow e_3}
    \and
    \inferrule*[right=\(\TirName{If}\)]
       {e_1 \rightarrow e'_1}
       {\Xif \; e_1 \; \Xthen \; e_2 \; \Xelse \; e_3
        \rightarrow \Xif \;
        e'_1 \; \Xthen \; e_2 \; \Xelse \; e_3}
  \end{mathpar}
  This means that the boolean expression \(e_1\) must be completely
  reduced \emph{before} the others. This is the usual way in
  programming languages, the rationale being to avoid unnecessary
  computations (evaluating \(e_2\) and~\(e_3\) is not necessary).

\end{frame}
